\scanpage{54}

en tous cas. De plus, bien entendu, pour ce qu'on veut en faire pour
R.R. , c'est bien la nouvelle notion d'immersion régulière qu\deleted{\unsure{e}}i fait
tout marcher.

Du point de vue pratique, il me semble évidemment trop tard pour
faire les changements désirables sur épreuves dans IV 16,19. On peut
cependant ajouter une remarque 19.1.8.,dans laquelle nous ferions l'aveu
\margin{\handwritten{\rule{0.25pt}{4.8cm}\hspace{2pt}\rule{0.25pt}{4.8cm}}}que \deleted{\unsure{la|ln}} la définition adoptée dans le texte de immersion régulière \deleted{n}s'est
\deleted{pas} avérée peu maniable dans le cas non noethérien, et qu'il y aurait
lieu de la modifier comme je viens de le dire, en\deleted{\unsure{r|é}} référant le lecteur
intéressé, \inserted{pour des détails} en attendant la réédition du Chap IV des EGA, au séminaire
SGA 6 (qui devrait sortir vers la fin de l'année 67 ou début 68). De
plus, à l'endroit du par.16 ou on introduit la définition inc\deleted{\unsure{i}}riminée,
nous pourrions insérer une note de bas de page, renvoyant le le cteur
à la remarque 19.1.8. concernant cette définition. (NB Dans le par.16,
c'est en fait la notion quasi-régulière qui sert vraiment ...).

Au moment de la réédition du Chap IV, il y aura lieu également de
revoir $0_{\mathrm{IV}}$ 15 dans le même esprit; je prévois que l'on pourra également
y remplacer partout les hypothèses ``suite régulière'', par l\deleted{a}'hypothèse
plus faible et plus maniable ``suite Koszul-régulière''. S'il en est
bien ainsi, il s'imposera alors, pour les suites également, de réserver
le mot le plus court à la notion qui manifestement est la plus serviable,
et il faudra envisaher de changer \deleted{\unsure{not|mot}} notre terminologie, en disant
``suite régulière'' au lieu de ``suite Koszul-régulière'', et en réservant
tout au plus un terme tel que ``suite strictement régulièr\deleted{\unsure{e|s}}e'' à la notion
éculée. Bien sûr, dans le cas noethérien, les trois variantes (régulière,
quasi; et Koko) sont équivalentes.
% archive sheet number: 53
